\subsection*{停在原处也可以满足安全条件}

上一节证明的安全要求限制错误能够发生什么：不产生同 \(S\) 的第二份扣款，不把独立身份误合并，不把原结果改写为另一份业务事实。数据库现在完全不可达，所有恢复者停止，原记录继续保留，这些要求仍可能成立。系统没有前进，却也没有增加错误状态。要使尚未形成的购买最终结束，还需逐项建立可用性、调度和持续恢复的条件。

进展讨论的是尚未结束的事情能否继续。例如 \(S\) 未形成时，最终会不会有一份终态；\(S\) 已成功时，客户端最终能否读回；客户端读回以后，又能否在本地可靠保存。这些目标依赖的条件各不相同。已经提交 P1 后，服务端业务完成目标已满足，客户端获知目标还可能因每次响应丢失而始终未满足。

Birrell 与 Nelson 在《Implementing Remote Procedure Calls》中把正常返回和通信异常后的语义分开：通信异常可能让调用方不知道远端执行了还是没有执行\cite{rpc}。本例保留 \(S\) 和持久结果，使暂时异常以后还能恢复。调用方遇到通信异常后继续保留 \(S\)，在重新获得执行或读取机会时，依据数据库终态确定结论。

\subsection*{从一个等待者推导完成条件}

设 A 已插入 S1，但在写订单以前停住，B 在唯一检查处等待。即使数据库仍在线，B 也无法自行把 A 的事务判为回滚。若 A 永远挂起且数据库永远保留连接与锁，B 可以一直等待。要保证继续，必须使占有资源的旧事务最终提交或被确定撤销，相关锁最终释放，后来的恢复得到运行机会。

实际进程终止实验满足了这个条件的一种具体情形：终止 A 导致其连接结束，数据库撤销未提交事务，等待者继续。服务中的超时、连接失效检测和受控取消也可帮助释放资源，但都需要数据库真正结束事务。应用单方面把内存标记改为可用，再跳过数据库占用，会让旧 A 恢复后与新 B 同时写入，破坏安全。

要使一项尚无终态的有效 \(S\) 最终完成，可以明确给出一组足够条件：保留 \(S\) 与原输入的恢复者不会永久停止；数据库最终提供足够长的可用时段；旧事务最终终结，锁不会永久被占用；请求与数据库执行不会一直被其他工作挤出；报价与余额判断能够在有限步骤内结束；许可及记录保留的时间足以让这一处理完成。在这些条件下，某次恢复可以取得 \(S\)，完成成功或拒绝分支并提交。

这里的调度条件通常称为公平性：一个持续有资格执行的恢复动作，最终能够得到调度机会。公平性是对无限执行历史的要求，有限次数观察到线程运行正常不足以证明它。也无需承诺每个竞争者都取得身份；同 \(S\) 只需一个尝试完成，其他尝试后来读取该终态即可。

\begin{figure}[htbp]
\centering
\begin{tikzpicture}
\node[booknode,text width=3.1cm] (make) at (1.7,2.2) {形成完整终态\\成功或业务拒绝};
\node[booknode,text width=3.1cm] (read) at (5.7,2.2) {读取原结果\\按完整 S 与原版本};
\node[bookresult,text width=3.1cm] (save) at (9.7,2.2) {客户端获知并保存\\完整交付\\本地持久化};
\draw[bookarrow] (make) -- (read);
\draw[bookarrow] (read) -- (save);
\node[booknote,text width=3.1cm] at (1.7,.4) {数据库可写\\旧事务终结释放锁\\恢复获得调度};
\node[booknote,text width=3.1cm] at (5.7,.4) {权威数据可读\\资格与记录仍有效\\解释版本可用};
\node[booknote,text width=3.1cm] at (9.7,.4) {响应完整到达\\客户端继续运行\\至少一次保存成功};
\end{tikzpicture}
\caption{完成终态、读取和客户端保存分别依赖后续机会}
\label{fig:ch06-progress-stages}
\end{figure}


成功扣款还多一项业务前提。若恢复最终得到执行机会时余额不足，协议应提交拒绝；持续重试无法把这个正确拒绝变成必定成功。对本章的进展目标，更准确的说法是最终得到一项完整业务结论。若产品另需保证某项购买一定成功，就须确保报价、资金与授权等成功条件在实际判断时满足。

\subsection*{已有终态以后还缺哪段机会}

\(S\) 已有完整终态时，恢复无需再次购买，只需能够定位并读回。此时前提变为权威数据库可读、完整 \(S\) 仍在恢复记录中、访问资格有效、旧输入与结果版本仍可解释、读请求能够得到调度。新连接即使经历等待，也应最终走到原结果。后台恢复者从自己的持久队列取得 \(S\)，可以在客户端进程退出后继续；只有内存里一个定时器，进程消失时便不再提供恢复机会。

客户端最终获知还要求至少一次完整结果能从服务到达它。数据库已经十次返回 P1，每次 HTTP 响应都在路上丢失，客户端仍可一直处于待确认。若响应到达后、本地保存以前客户端退出，也可能在重启后重新查询；保留 \(S\) 使这仍安全，但要承诺最终保存，还需一次本地持久化成功的机会。服务器发送日志记录服务端已经执行发送步骤；客户端完整收到响应并成功保存本地结果以后，获知与持久保存这两个目标才都得到满足。

重试间隔影响资源与机会。所有客户端同时立即重发，会增加锁等待与数据库负担。退避可以让每次失败后逐渐延长等待，并加少量随机偏移，避免恢复请求总在同一时刻集中；它是调度负载的办法，不改变同 \(S\) 的事务规则。每次尝试可以有有限等待上限，超时后返回 503 暂时不可用，客户端保留身份，再等下一机会。

有限次重试与最终完成也有明显差别。客户端只试三次，恰好三次都在数据库不可用期间结束，第四次本可成功却从未发送，进展条件就不成立。另一方面，把单次等待设为无限，也无法保证数据库或持锁事务最终恢复。恢复须持续到足够的可用时段，旧事务在此期间终结并释放锁，后续恢复获得实际运行机会。

\subsection*{固定截止时间怎样限制最终获知}

E1 采用固定截止时刻。设调用方断网直到截止之后，旧入口返回 \code{410 operation\_expired}。410 表示这份恢复资格已经结束，原数据库可能保存成功，也可能尚未形成终态；它不证明扣款失败。若没有独立业务事实核查，就自动改 K2 重买，先前已经成功的 K1 仍可与 K2 形成两笔购买。

因此，对有限许可给出的进展保证只能覆盖有足够恢复窗口的历史。以 \(T\) 表示 E1 截止时刻，以 \(t\) 表示恢复开始时刻；可用时间为 \(T-t\)。要让原入口完成读取、交付与本地保存，相关步骤必须在允许规则内获得足够机会。这里没有规定网络延迟的确定上限，所以仅凭 \(t<T\) 无法推出必在 \(T\) 前获知；临近截止的一次有效发送仍可能错过最终交付。

截止也不能成为清理未结束事务的直接理由。第三章的关闭边界要求先阻止新受理，保留在途执行的身份与结果直至安全终结，随后才清理；否则，前面归纳中依赖的记录保留条件就会被破坏。

本章有界模型把许可有效、记录不回收列为前提，实际进程实验也在这一范围内运行。完整回收协议需要另行验证。发生长期断网或许可到期以后，客户端继续保留原身份及可核对的购买信息，经账户内业务事实查询或受控处理确认原结论；需要另一项购买时，再由用户明确确认并分配新身份。
\FloatBarrier
